\chapter{社会认知}

\section{社会认知领域}

社会认知（Social Cognition）维护其语义依赖特定个人、二人关系、群体、频道或社区的长期认知。领域至少包含关系认知图和语言惯例两类一等对象族。

普通世界知识和通用语言知识不因涉及人物而进入社会认知。只有在去除具体社会作用域后语义会发生变化的长期认知，才由社会认知领域拥有。

\section{关系认知图}

关系认知图保存某个主体对社会世界的长期认识。每项关系认知都属于一个拥有者主体，并连接来源实体、关系类型和目标实体。关系默认具有方向性。

关系认知至少包含以下语义维度：

\begin{itemize}
\item 拥有该认知的主体；
\item 关系的来源实体和目标实体；
\item 关系类型；
\item 在该关系类型具有程度语义时的可选关系程度；
\item 主体对该关系判断的认识状态；
\item 支持与反驳证据、有效时间和修订历史。
\end{itemize}

关系程度与认识确定程度分别表示。主体可以低确定度地认为两人关系非常亲密，也可以高确定度地认为两人只有普通联系。不同关系类型的程度值没有默认的跨类型算术意义。

关系图中的路径用于检索、相关对象发现和假设形成。关系不具有默认传递性；由多跳路径推导出的新关系首先形成假设，经过证据搜索和社会认知领域核验后才能提交。

\section{语言惯例}

语言惯例（LanguageConvention）保存特定社会作用域内形成的词义、称呼、黑话、简称和语用约定。每项语言惯例至少包含表达形式、社会作用域、语义或指代、可选语用角色、认识状态、证据、有效时间和修订。

社会作用域可以是个人、二人关系、群体、社区或频道。查询时依据当前参与者和场景选择适用的惯例；不同作用域可以同时保留彼此不同的含义。

显式解释、反复一致使用、成功指代、确认和修复过程都可以形成语言惯例的证据。单次上下文猜测只能形成假设。主体自己在学习后重复使用某个表达，不构成该惯例新的独立外部证据；独立支持来自新的外部观察、对方继续使用、正确理解或显式确认。

\section{实体身份}

社会认知引用已经绑定的实体身份。昵称、显示名称和向量相似度不直接合并人物身份。跨平台身份关联保存独立来源和证据；平台好友关系、群角色和账号信息可以作为来源事实，再由社会认知领域形成长期关系认知。
